
    %%%%%%%%%%%%%%%%%%%%%%%%%%%%%%%%%%%%%%%%%%%%%%%%%%%%%%%%%%%%%%%%%%%%%%%%%%%%%%%%
    %%% Classe do documento
    \documentclass[12pt, addpoints]{exam}
    \UseRawInputEncoding

    %%%%%%%%%%%%%%%%%%%%%%%%%%%%%%%%%%%%%%%%%%%%%%%%%%%%%%%%%%%%%%%%%%%%%%%%%%%%%%%%
    % preambulo.tex
    %
    % Arquivo de configuração de packages para uso o LaTeX, conforme minhas
    % preferências e modelos pessoais.
    %
    % Para maiores informações, visite:
    %    https://github.com/abrantesasf/latex
    %
    % NÃO ALTERE SE NÃO SOUBER O QUE ESTÁ FAZENDO!
    %%%%%%%%%%%%%%%%%%%%%%%%%%%%%%%%%%%%%%%%%%%%%%%%%%%%%%%%%%%%%%%%%%%%%%%%%%%%%%%%


    %%%%%%%%%%%%%%%%%%%%%%%%%%%%%%%%%%%%%%%%%%%%%%%%%%%%%%%%%%%%%%%%%%%%%%%%%%%%%%%%
    %%% Estruturas de controle:
    \input{static/utils/controles}

    %%%%%%%%%%%%%%%%%%%%%%%%%%%%%%%%%%%%%%%%%%%%%%%%%%%%%%%%%%%%%%%%%%%%%%%%%%%%%%%%
    %%% Compilação condicional em PDF:
    \input{static/utils/pdf}

    %%%%%%%%%%%%%%%%%%%%%%%%%%%%%%%%%%%%%%%%%%%%%%%%%%%%%%%%%%%%%%%%%%%%%%%%%%%%%%%%
    %%% Configurações de layout da página:
    \input{static/utils/layout}

    %%%%%%%%%%%%%%%%%%%%%%%%%%%%%%%%%%%%%%%%%%%%%%%%%%%%%%%%%%%%%%%%%%%%%%%%%%%%%%%%
    %%% Configurações para autores:
    \input{static/utils/autores}

    %%%%%%%%%%%%%%%%%%%%%%%%%%%%%%%%%%%%%%%%%%%%%%%%%%%%%%%%%%%%%%%%%%%%%%%%%%%%%%%%
    %%% Cabeçalho e rodapé:
    %%% Atenção! É MUITO DIFÍCIL ajustar apropriadamente os running headers
    %%% e running footers da primeira vez. Você pode confiar no LaTeX e deixar que
    %%% ele faça o ajuste automaticamente ou pode quebrar a cabeça e
    %%% tentar melhorar algo que o LaTeX já faz muito bem, mesmo que não fique
    %%% exatamente do jeito que você quer. O que você prefere? Se quiser ajustar
    %%% manualmente, descomente a linha abaixo e faça as configurações no arquivo
    %%% "utils/headings.tex".
    %\input{static/utils/headings}

    %%%%%%%%%%%%%%%%%%%%%%%%%%%%%%%%%%%%%%%%%%%%%%%%%%%%%%%%%%%%%%%%%%%%%%%%%%%%%%%%
    %%% Configuração para as fontes do documento:
    %%% Se você quiser usar fontes próprias ou do sistema, precisará ajustar as
    %%% configurações no arquivo "utils/fontes.tex".
    %%% ATENÇÃO: por padrão eu utilizo XeLaTeX com fontes proprietárias projetadas
    %%% por Matthew Butterick, adquiridas em https://mbtype.com, que não podem ser
    %%% distribuídas devido à licença de utilização. Se você usar XeLaTeX, você
    %%% DEVERÁ OBRIGATORIAMENTE ajustar as fontes no arquivo "utils/fontes.tex".
    \input{static/utils/fontes}

    %%%%%%%%%%%%%%%%%%%%%%%%%%%%%%%%%%%%%%%%%%%%%%%%%%%%%%%%%%%%%%%%%%%%%%%%%%%%%%%%
    %%% Sumário:
    \input{static/utils/toc}

    %%%%%%%%%%%%%%%%%%%%%%%%%%%%%%%%%%%%%%%%%%%%%%%%%%%%%%%%%%%%%%%%%%%%%%%%%%%%%%%%
    %%% Matemática:
    \input{static/utils/matematica}

    %%%%%%%%%%%%%%%%%%%%%%%%%%%%%%%%%%%%%%%%%%%%%%%%%%%%%%%%%%%%%%%%%%%%%%%%%%%%%%%%
    %%% Notas de rodapé:
    \input{static/utils/footnotes}

    %%%%%%%%%%%%%%%%%%%%%%%%%%%%%%%%%%%%%%%%%%%%%%%%%%%%%%%%%%%%%%%%%%%%%%%%%%%%%%%%
    %%% Referências cruzadas, links, citações:
    \input{static/utils/crossrefs}

    %%%%%%%%%%%%%%%%%%%%%%%%%%%%%%%%%%%%%%%%%%%%%%%%%%%%%%%%%%%%%%%%%%%%%%%%%%%%%%%%
    %%% Configurações para os hiperlinks em PDF:
    \input{static/utils/pdfconfig}

    %%%%%%%%%%%%%%%%%%%%%%%%%%%%%%%%%%%%%%%%%%%%%%%%%%%%%%%%%%%%%%%%%%%%%%%%%%%%%%%%
    %%% Glossário e índice remissivo:
    \input{static/utils/glossindex}

    %%%%%%%%%%%%%%%%%%%%%%%%%%%%%%%%%%%%%%%%%%%%%%%%%%%%%%%%%%%%%%%%%%%%%%%%%%%%%%%%
    %%% Referências bibliográficas:
    \input{static/utils/bibliografia}

    %%%%%%%%%%%%%%%%%%%%%%%%%%%%%%%%%%%%%%%%%%%%%%%%%%%%%%%%%%%%%%%%%%%%%%%%%%%%%%%%
    %%% Cores:
    \input{static/utils/cores}

    %%%%%%%%%%%%%%%%%%%%%%%%%%%%%%%%%%%%%%%%%%%%%%%%%%%%%%%%%%%%%%%%%%%%%%%%%%%%%%%%
    %%% Computação:
    \input{static/utils/computacao}

    %%%%%%%%%%%%%%%%%%%%%%%%%%%%%%%%%%%%%%%%%%%%%%%%%%%%%%%%%%%%%%%%%%%%%%%%%%%%%%%%
    %%% Gráficos:
    \input{static/utils/graficos}

    %%%%%%%%%%%%%%%%%%%%%%%%%%%%%%%%%%%%%%%%%%%%%%%%%%%%%%%%%%%%%%%%%%%%%%%%%%%%%%%%
    %%% Ícones:
    \input{static/utils/icones}

    %%%%%%%%%%%%%%%%%%%%%%%%%%%%%%%%%%%%%%%%%%%%%%%%%%%%%%%%%%%%%%%%%%%%%%%%%%%%%%%%
    %%% Blocos:
    \input{static/utils/blocos}

    %%%%%%%%%%%%%%%%%%%%%%%%%%%%%%%%%%%%%%%%%%%%%%%%%%%%%%%%%%%%%%%%%%%%%%%%%%%%%%%%
    %%% Boxes coloridos:
    \input{static/utils/colorbox}

    %%%%%%%%%%%%%%%%%%%%%%%%%%%%%%%%%%%%%%%%%%%%%%%%%%%%%%%%%%%%%%%%%%%%%%%%%%%%%%%%
    %%% Tabelas:
    \input{static/utils/tabelas}

    %%%%%%%%%%%%%%%%%%%%%%%%%%%%%%%%%%%%%%%%%%%%%%%%%%%%%%%%%%%%%%%%%%%%%%%%%%%%%%%%
    %%% Ambientes de listas:
    \input{static/utils/listas}

    %%%%%%%%%%%%%%%%%%%%%%%%%%%%%%%%%%%%%%%%%%%%%%%%%%%%%%%%%%%%%%%%%%%%%%%%%%%%%%%%
    %%% Ambientes floats e similares:
    \input{static/utils/floats}

    %%%%%%%%%%%%%%%%%%%%%%%%%%%%%%%%%%%%%%%%%%%%%%%%%%%%%%%%%%%%%%%%%%%%%%%%%%%%%%%%
    %%% Ajustes para a classe EXAM:
    \input{static/utils/exam}

    %%%%%%%%%%%%%%%%%%%%%%%%%%%%%%%%%%%%%%%%%%%%%%%%%%%%%%%%%%%%%%%%%%%%%%%%%%%%%%%%
    %%% Ajustes para textos:
    \input{static/utils/texto}

    %%%%%%%%%%%%%%%%%%%%%%%%%%%%%%%%%%%%%%%%%%%%%%%%%%%%%%%%%%%%%%%%%%%%%%%%%%%%%%%%
    %%% Ambientes, aliases, comandos e customizações em geral para serem
    %%% utilizadas nos documentos. Defina aqui qualquer customização específica
    %%% para seus documentos!
    \input{static/utils/customizacoes}

    %%%%%%%%%%%%%%%%%%%%%%%%%%%%%%%%%%%%%%%%%%%%%%%%%%%%%%%%%%%%%%%%%%%%%%%%%%%%%%%%
    %%% Ajuste do layout, espaçamento de linhas e etc.:
    \geometry{a4paper, portrait, left=2.0cm, right=2.0cm, top=2.0cm, bottom=2.0cm}
    \singlespacing

    %%%%%%%%%%%%%%%%%%%%%%%%%%%%%%%%%%%%%%%%%%%%%%%%%%%%%%%%%%%%%%%%%%%%%%%%%%%%%%%%
    %%% Configurações de cabeçalho e rodapé (não use fancyhdr pois ocorre conflito):
    \pagestyle{headandfoot}
    \runningheadrule
    \firstpageheader{}{}{}
    \runningheader{Sem disciplina}{}{Avaliação}
    \firstpagefooter{}{}{Página \thepage\ de \numpages}
    \runningfooter{}{}{Página \thepage\ de \numpages}
    \runningfootrule

    %%%%%%%%%%%%%%%%%%%%%%%%%%%%%%%%%%%%%%%%%%%%%%%%%%%%%%%%%%%%%%%%%%%%%%%%%%%%%%%%
    %%% Configurações para as propriedades do PDF:
    \hypersetup{
    hidelinks,           % Comente para web, descomente para impressão
    colorlinks=true,      % True para web, False para impressão
    pdftitle=Pedro: Avaliação,
    pdfauthor=Matheus Endlich Silveira,
    pdfsubject=Sem disciplina,
    pdfkeywords=['Sem tag'],
    pdfinfo={
        CreationDate={}, % Ex.: D:AAAAMMDDHH24MISS
        ModDate={}       % Ex.: D:AAAAMMDDHH24MISS
    }
    }
    \input{static/utils/pdfconfig}

    %%%%%%%%%%%%%%%%%%%%%%%%%%%%%%%%%%%%%%%%%%%%%%%%%%%%%%%%%%%%%%%%%%%%%%%%%%%%%%%%
    %%% Começa o documento
    \begin{document}

    %%%%%%%%%%%%%%%%%%%%%%%%%%%%%%%%%%%%%%%%%%%%%%%%%%%%%%%%%%%%%%%%%%%%%%%%%%%%%%%%
    %%% Folha de instruções padronizadas da UVV para a prova
    \pdfbookmark[1]{Instruções}{instrucoes}
    \begin{figure}[H]
        \begin{center}
            \includegraphics[scale=0.3]{{static/imgs/logo-uvv.png}}
        \end{center}	
    \end{figure}

    \vspace{-1cm}

    \begin{table}[H]
        \centering
        \begin{tabular}{|llll|}
            \hline
            \multicolumn{2}{|l|}{\textbf{Disciplina:} Sem disciplina} & 
            \multicolumn{1}{l|}{\multirow{3}{*}{\textbf{\makecell{Nota:\\ \,\\ \,}}}} & 
            \multirow{3}{*}{\textbf{\makecell{Coordenador:\\ \,\\ \,}}} \\ 
            \cline{1-2}
            \multicolumn{2}{|l|}{\textbf{Professor:} Matheus Endlich Silveira \hspace{4.72cm} } & 
            \multicolumn{1}{l|}{} & \\ 
            \cline{1-2}
            \multicolumn{2}{|l|}{\textbf{Aluno:}} & 
            \multicolumn{1}{l|}{} & \\ 
            \hline
            \multicolumn{1}{|l|}{\textbf{Turma:} \hspace{6.3cm} } & 
            \multicolumn{1}{l|}{\textbf{Semestre:}} & 
            \multicolumn{2}{l|}{\textbf{Valor:} 10 pontos} \\ 
            \hline
            \multicolumn{1}{|l|}{\textbf{Data:}} & 
            \multicolumn{3}{l|}{\textbf{Avaliação:}} \\ 
            \hline
        \end{tabular}
    \end{table}

    \begin{center}
        \fbox{\setlength\fboxsep{0.3cm} \fbox{\parbox{15.4cm}{
            \begin{center}
                \textbf{INSTRUÇÕES PARA A REALIZAÇÃO DESTA AVALIAÇÃO:}
            \end{center}
            \begin{itemize}[leftmargin=*]
                \item Esta é a primeira avaliação da disciplina Sem disciplina, 
                    e engloba o conteúdo referente Sem tag.
                \item Esta avaliação contém 12 questões objetivas, \textbf{todas obrigatórias}.
                \item \textbf{Leia atentamente cada questão!} As questões objetivas terão uma,
                    e apenas uma única, resposta a ser marcada.
                \item \textbf{Desligue o celular} antes de começar. A avaliação é
                    \textbf{individual} e \textbf{sem consulta}.
                \item As questões podem ser respondidas, na prova, com lápis ou caneta. Ao final
                    da prova \textbf{{VOCÊ DEVERÁ PREENCHER O GABARITO COM CANETA
                    DE COR PRETA}} \textbf{\underline{OU AZUL}} para que seja feita a correção
                    automatizada através da leitura do padrão de respostas marcado no
                    gabarito.
                \item \textbf{CUIDADO ao preencher o gabarito} pois eles são \textbf{INDIVIDUAIS
                    e NOMEADOS} (cada estudante tem um gabarito próprio específico, com um
                    código de identificação único). \textbf{Questões rasuradas são anuladas}
                    pelo software de leitura do gabarito.
                \item Ao final da prova, \textbf{DEVOLVA A PROVA E O GABARITO} para o professor.
                \item Siga todas as normas de \textbf{Integridade Acadêmica} da disciplina.
                    Alunos que forem flagrados com qualquer espécie de ``cola'' ou trocando
                    informações com outros alunos terão suas avaliações recolhidas, as notas
                    zeradas, e a situação será encaminhada para a coordenação para a aplicação
                    das penalidades previstas pela universidade.
                \item Com 90 minutos de avaliação você terá tempo de até 8,min por questão,
                    em média.
                \item Boa avaliação!
            \end{itemize}
            \vspace{2cm}
        }}}
    \end{center}
    \newpage

    %%%%%%%%%%%%%%%%%%%%%%%%%%%%%%%%%%%%%%%%%%%%%%%%%%%%%%%%%%%%%%%%%%%%%%%%%%%%%%%%
    %%% Inicia o bloco de questões:
    \begin{questions}

    \newpage
    \question

Um compilador detecta:
\begin{choices}
\choice erros aritméticos
\choice erros de sintaxe do programa 
\choice erros que podem ocorrer durante a execução do programa
\choice todos os erros citados acima
\choice erros nos resultados gerados pelo programa
\end{choices}

\question

(ENADE 2017) No modelo relacional de dados, uma junção entre tabelas cria uma

outra tabela derivada de duas ou mais tabelas de acordo com os critérios

especificados para a junção, de modo similar às regras da teoria dos

conjuntos. Nos conjuntos a seguir, as áreas hachuradas representam os resultados

das consultas SQL em que ``id' é uma chave primária e ``fk' é uma chave

estrangeira:

\begin{figure}[H]

  \centering

  \includegraphics[width=0.5\linewidth]{static/imgs/defaul.png}

\end{figure}

Assinale a opção em que a declaração SQL é equivalente ao conjunto apresentado a

seguir:

\begin{figure}[H]

  \centering

  \includegraphics[width=0.5\linewidth]{static/imgs/defaul.png}

\end{figure}


\begin{choices}
\choice \begin{verbatim}SELECT *

FROM tabela1 t1

WHERE NOT EXISTS (SELECT 1

                  FROM tabela1 t2

                  WHERE t2.id = tk.fk);

\end{verbatim}


\choice \begin{verbatim}SELECT *

FROM tabela1 t1 WHERE EXISTS (SELECT 1

                              FROM tabela2 t2

                              WHERE t2.id = t1.fk);

\end{verbatim}


\choice \begin{verbatim}SELECT *

FROM tabela1 t1

INNER JOIN tabela2 t2 ON (t1.id = t2.fk);

\end{verbatim}


\choice \begin{verbatim}SELECT *

FROM tabela1 t1

RIGHT JOIN tabela2 t2 ON (t1.id = t2.fk)

WHERE t1.id IS NULL;

\end{verbatim}


\choice \begin{verbatim}SELECT *

FROM tabela1 t1

RIGHT JOIN tabela2 t2 ON (t1.id = t2.fk)

WHERE t2.fk <> NULL;

\end{verbatim}


\end{choices}

\question

Filtro da mediana é:
\begin{choices}
\choice Indicado para detectar formas específicas em imagens.
\choice Indicado para detectar tonalidades específicas em uma imagem.
\choice Nenhuma das respostas acima.
\choice Indicado para detectar bordas em imagens.
\choice Indicado para atenuar ruído com preservação de bordas (i.e., rápidas transições de nível em imagens).
\end{choices}

\question

Uma prova de vestibular foi elaborada com 25 questões de múltipla escolha com 5 alternativas. O número de candidatos presentes à prova foi 63127. Considere a afirmação: Pelo menos 2 candidatos responderam de modo idêntico as \( k \) primeiras questões da prova. Qual é o maior valor de \( k \) para o qual podemos garantir que a afirmação é verdadeira.


\begin{choices}
\choice 6
\choice 9
\choice 10
\choice 7
\choice 8
\end{choices}

\question

Em um banco de dados relacional, o que é uma chave primária?


\begin{choices}
\choice Nenhuma das respostas acima
\choice É uma validação que evita o cadastro de dados errados
\choice É denominada por FK
\choice É um atributo (ou um conjunto de atributos) que permite identificar única e exclusivamente cada registro da tabela
\choice É qualquer coluna ou conjunto de colunas que podem servir como índice
\end{choices}

\question

Considere as seguintes afirmações sobre redes neurais artificiais: \begin{enumerate}[label=\Roman*.] \item Um perceptron elementar só computa funções linearmente separáveis. \item Não aceitam valores numéricos como entrada. \item O "conhecimento" é representado principalmente através do peso das conexões. \end{enumerate} São corretas:
\begin{choices}
\choice I, II e III
\choice Apenas III
\choice Apenas I e II 
\choice Apenas II e III
\choice Apenas I e III
\end{choices}

\question

Considerando a rede de Petri abaixo, quais das alternativas são verdadeiras?



\begin{itemize}

    \item[I)] O lugar \( A \) está habilitado a disparar.

    \item[II)] Apenas a transição \( T1 \) está habilitada a disparar.

    \item[III)] A sequência de transições \( (T1, T2, T3, T2) \) pode ser disparada, nessa ordem.

    \item[IV)] A transição \( T4 \) nunca poderá ser disparada.

\end{itemize}



\begin{figure}[H]

    \centering

    \includegraphics[width=0.5\linewidth]{static/imgs/defaul.png}

    \caption{Questão 49}

    \label{fig:enter-label}

\end{figure}
\begin{choices}
\choice Apenas as alternativas II e III.
\choice Apenas as alternativas I e III.
\choice Apenas as alternativas III, IV.
\choice Todas as alternativas.
\choice Apenas as alternativas II, III e IV.
\end{choices}

\question

Sobre a hierarquia de Chomsky podemos afirmar que:
\begin{choices}
\choice Há linguagens que não são nem livres de contexto nem sensíveis a contexto
\choice As linguagens livres de contexto e as linguagens sensíveis a contexto se excluem
\choice Uma linguagem que não é regular é livre de contexto
\choice Uma linguagem que é recursivamente enumerável não pode ser uma linguagem regular
\choice As linguagens reconhecidas por autômatos a pilha são as linguagens regulares
\end{choices}

\question

Quando dizemos que um banco de dados é em \ingles{cloud} e centralizado, estamos

nos referindo, respectivamente, à localização dos dados e da

infraestrutura. Verdadeiro ou falso?
\begin{choices}
\choice Verdadeiro
\choice Não é possível determinar com os dados apresentados.
\choice Falso
\end{choices}

\question

Uma definição lógica, abstrata e auto-contida dos objetos que modelam a

estrutura de armazenamento de dados, dos operadores que modelam o comportamento

e dos demais conceitos que modelam a integridade, e que, juntos, constituem a

máquina abstrata com a qual os usuários interagem é chamada de:


\begin{choices}
\choice Banco de dados
\choice Modelo de dados
\choice Modelo entidade-relacionamento
\choice Modelo relacional
\choice Nenhuma das respostas acima
\end{choices}

\question

Em qual das situações abaixo um sistema de Raciocínio Baseado em Casos não deve ser utilizado?


\begin{choices}
\choice Quando especialistas conversam sobre seus domínios dando exemplos.
\choice Quando for fácil a obtenção de regras
\choice Quando a experiência for tão valiosa quanto o conhecimento em livros texto.
\choice Quando as regras utilizadas apresentam um grande número de exceções.
\choice Em aplicações de diagnóstico médico.
\end{choices}

\question

Considere as seguintes afirmativas:

\begin{itemize}

    \item O modelo matemático de uma lista é a sequência linear de itens, cuja principal propriedade estrutural é a posição relativa dos elementos dentro da sequência.

    \item A fila e a pilha são consideradas casos especiais da lista.

    \item Numa fila a inserção e a retirada são feitas no mesmo extremo.

    \item Numa lista a inserção e a retirada podem ser feitas em qualquer posição.

    \item Numa pilha apenas a inserção pode ser feita em qualquer posição.

\end{itemize}



Quais são as afirmativas verdadeiras?
\begin{choices}
\choice somente II, III e IV
\choice somente I, II e IV
\choice somente I e III
\choice somente II, IV e V
\choice todas
\end{choices}



    %%%%%%%%%%%%%%%%%%%%%%%%%%%%%%%%%%%%%%%%%%%%%%%%%%%%%%%%%%%%%%%%%%%%%%%%%%%%%%%%
    %%% Finaliza o bloco de questões:
    \end{questions}
    \end{document}
    